% Palimpsest — Linux LaTeX header for `pal pdf` inside the container.
% Same typography as the macOS default, but with fonts that ship in the image
% (texlive-fonts-recommended + fonts-texgyre + fonts-noto + fonts-noto-cjk), so
% XeLaTeX finds every glyph. The container points at this via PAL_PDF_HEADER.

\usepackage{fancyhdr}

% --- Prose typography (book-like) ---------------------------------------
\usepackage{microtype}
\usepackage{setspace}\doublespacing
\setlength{\parindent}{0pt}
\setlength{\parskip}{0.7\baselineskip plus 2pt}
\raggedbottom
\frenchspacing

% --- Fonts (all present in the image) -----------------------------------
% TeX Gyre Termes is a Times-like transitional serif; TeX Gyre Cursor is the
% Courier-like mono. Both come with texlive-fonts-recommended.
\setmainfont{TeX Gyre Termes}
\setmonofont{TeX Gyre Cursor}[Scale=0.85]

% --- Multilingual glyph fallback ----------------------------------------
% Auto-switch fonts per Unicode block via XeTeX interchar classes (same
% technique as the macOS header). Noto covers Greek, Hebrew, arrows, and CJK.
\XeTeXinterchartokenstate=1
\newcount\ebloop

\def\ebrange#1#2#3{%
  \ebloop=#2
  \loop \ifnum\ebloop<#3
    \XeTeXcharclass\ebloop=#1
    \advance\ebloop by 1
  \repeat}
\def\ebbind#1#2{%
  \XeTeXinterchartoks 0 #1 = {\bgroup#2}%
  \XeTeXinterchartoks 255 #1 = {\bgroup#2}%
  \XeTeXinterchartoks 4095 #1 = {\bgroup#2}%
  \XeTeXinterchartoks #1 0 = {\egroup}%
  \XeTeXinterchartoks #1 255 = {\egroup}%
  \XeTeXinterchartoks #1 4095 = {\egroup}}

% Greek (U+0370–03FF) + Greek Extended (U+1F00–1FFF).
\newfontfamily\greekfont{Noto Serif}
\newXeTeXintercharclass\greekclass
\ebrange\greekclass{"0370}{"0400}
\ebrange\greekclass{"1F00}{"2000}
\ebbind\greekclass\greekfont

% Hebrew (U+0590–05FF).
\newfontfamily\hebrewfont{Noto Sans Hebrew}
\newXeTeXintercharclass\hebrewclass
\ebrange\hebrewclass{"0590}{"0600}
\ebbind\hebrewclass\hebrewfont

% Arrows (U+2190–21FF) + letterlike aleph (transfinite orders).
\newfontfamily\arrowfont{Noto Sans Symbols}
\newXeTeXintercharclass\arrowclass
\ebrange\arrowclass{"2190}{"2200}
\ebrange\arrowclass{"2135}{"2139}
\ebbind\arrowclass\arrowfont

% CJK Unified Ideographs (U+4E00–9FFF).
\newfontfamily\cjkfont{Noto Sans CJK SC}
\newXeTeXintercharclass\cjkclass
\ebrange\cjkclass{"4E00}{"A000}
\ebbind\cjkclass\cjkfont
